\section{Synchronization}
A \textcolor{Blue}{\textbf{race condition}} occurs when concurrent accesses to a shared resource yield non-deterministic outcomes due to interleaving (e.g., two threads doing \texttt{count++}).
A \textcolor{Blue}{\textbf{critical section}} is code that accesses such shared state.

\paragraph{Critical-section requirements}
\begin{itemize}
  \item \textcolor{Bittersweet}{Mutual exclusion}: at most one process inside at a time.
  \item \textcolor{Bittersweet}{Progress}: if no one is inside, an entrant must eventually be selected.
  \item \textcolor{Bittersweet}{Bounded waiting}: a finite bound on entries by others while a requester waits.
\end{itemize}

\subsection{Software solution: Peterson (2-process)}
\begin{verbatim}
flag[i] = true; turn = j;
while (flag[j] && turn == j) ; /* spin */
/* critical section */
flag[i] = false;
\end{verbatim}
Satisfies all three requirements; assumes atomic single-word writes and no compiler reordering.

\subsection{Hardware atomics}
\begin{itemize}
  \item \textcolor{Bittersweet}{Test-and-set} \texttt{TAS(\&lock)}: atomically reads old value, sets to \texttt{true}, returns old. Spinlock: \texttt{while (TAS(\&lock)) ;}.
  \item \textcolor{Bittersweet}{Compare-and-swap} \texttt{CAS(\&m, exp, new)}: if \texttt{*m==exp}, set to \texttt{new}, return true; else return false. Enables lock-free algorithms.
\end{itemize}
\textcolor{Bittersweet}{Spinlock vs blocking}: spin is fine for very short critical sections; blocking (semaphore) is better when contention is long because spinning wastes CPU.

\subsection{Semaphores}
A \textcolor{Blue}{semaphore} is an integer with two atomic operations:
\begin{itemize}
  \item \texttt{\textcolor{Bittersweet}{wait}(s)} (P/down): if $s>0$ then $s\!-\!-$, else block until \texttt{signal}.
  \item \texttt{\textcolor{Bittersweet}{signal}(s)} (V/up): $s\!+\!+$ and wake one waiter.
\end{itemize}
\textcolor{Bittersweet}{Binary} (saturates at 1; subsequent \texttt{signal} is a no-op) — equivalent to a mutex any thread can release. \textcolor{Bittersweet}{Counting} (unbounded; general resources). \texttt{sem\_init}, \texttt{sem\_wait}, \texttt{sem\_post}.

\subsection{Monitors}
A \textcolor{Blue}{monitor} is a high-level construct: only one process executes inside at a time. \textcolor{Bittersweet}{Condition variables} \texttt{cv} support \texttt{\textcolor{Bittersweet}{wait}(cv)} (release lock, sleep, re-acquire), \texttt{\textcolor{Bittersweet}{signal}(cv)}, \texttt{\textcolor{Bittersweet}{broadcast}(cv)}.
\begin{itemize}
  \item \textcolor{Bittersweet}{Hoare semantics}: signaled thread runs immediately; signaler resumes after.
  \item \textcolor{Bittersweet}{Mesa semantics}: signaled thread becomes Ready; signaler keeps the lock. Predicate may have changed by wake-up — \emph{must} re-check with \texttt{while}, not \texttt{if}. Reason: between signal and wakeup, another consumer can drain the buffer; on wakeup, with \texttt{if} the woken thread proceeds on a now-false predicate.
\end{itemize}

\subsection{Classic problems}

\paragraph{Bounded buffer (size $N$)}
Semaphores \texttt{mutex}=1, \texttt{empty}=$N$, \texttt{full}=0.
\begin{tabular}{@{}lll@{}}
\toprule
\# & producer & consumer \\
\midrule
init & \multicolumn{2}{l}{\texttt{mutex=1, empty=N, full=0}} \\
1 & \texttt{wait(empty)} & \texttt{wait(full)} \\
2 & \texttt{wait(mutex)} & \texttt{wait(mutex)} \\
3 & \texttt{insert} & \texttt{remove} \\
4 & \texttt{signal(mutex)} & \texttt{signal(mutex)} \\
5 & \texttt{signal(full)} & \texttt{signal(empty)} \\
\bottomrule
\end{tabular}
\textcolor{ForestGreen}{pthread / monitor form (Mesa)}:
\begin{verbatim}
pthread_mutex_lock(&m);
while (count == N)
    pthread_cond_wait(&notFull, &m);
buf[in] = x; in = (in+1)%N; count++;
pthread_cond_signal(&notEmpty);
pthread_mutex_unlock(&m);
\end{verbatim}

\paragraph{Readers-writers}
Multiple readers may hold simultaneously; writers exclusive.
\begin{itemize}
  \item \textcolor{ForestGreen}{Reader-preference}: readers count protected by \texttt{mutex}; first reader \texttt{wait}s on \texttt{rw}, last \texttt{signal}s. Writers may starve.
  \item \textcolor{ForestGreen}{Writer-preference}: pending writer blocks new readers.
\end{itemize}
\textcolor{ForestGreen}{Reader-preference (semaphore form)}:
\begin{verbatim}
reader: wait(mutex); rc++; if(rc==1) wait(rw); signal(mutex);
        /*read*/; wait(mutex); rc--; if(rc==0) signal(rw); signal(mutex);
writer: wait(rw); /*write*/; signal(rw);
\end{verbatim}

\paragraph{Dining philosophers}
$n$ philosophers, $n$ forks; each needs both adjacent forks. Naive symmetric solution \textcolor{Bittersweet}{deadlocks}. Fixes: pick up both atomically; allow at most $n-1$ to try; impose fork ordering (resource hierarchy); asymmetric (odd take left first, even take right first).

\paragraph{Priority inversion}
A high-priority task is blocked on a lock held by a low-priority task; an unrelated medium-priority task preempts the low-priority holder, indirectly delaying the high-priority task.
\textcolor{Bittersweet}{Priority inheritance}: temporarily raise the holder's priority to the highest waiter's priority until it releases.
\textcolor{Bittersweet}{Priority ceiling}: every lock has a ceiling = max priority of any task that may use it; holder runs at the ceiling.

\subsection{Deadlock}
Four necessary conditions: \textcolor{Bittersweet}{mutual exclusion}, \textcolor{Bittersweet}{hold-and-wait}, \textcolor{Bittersweet}{no preemption}, \textcolor{Bittersweet}{circular wait}.
Strategies:
\begin{itemize}
  \item \textcolor{Bittersweet}{Prevention}: negate one condition (e.g., total resource ordering breaks circular wait).
  \item \textcolor{Bittersweet}{Avoidance}: Banker's algorithm — see below.
  \item \textcolor{Bittersweet}{Detection \& recovery}: cycle search in resource-allocation graph; abort or rollback.
  \item \textcolor{Bittersweet}{Ostrich}: ignore (common when deadlock is rare).
\end{itemize}
\paragraph{Banker's algorithm}
Maintain Available, Max, Allocation, Need=Max$-$Allocation. Loop: find $i$ with $\text{Need}_i \le \text{Work}$; $\text{Work} \mathrel{+}= \text{Allocation}_i$; mark $i$ done; repeat. Safe iff every process gets marked.
Example: Available $=(3,3,2)$;
\begin{tabular}{@{}llll@{}}
\toprule
Proc & Max & Alloc & Need \\
\midrule
P0 & (7,5,3) & (0,1,0) & (7,4,3) \\
P1 & (3,2,2) & (2,0,0) & (1,2,2) \\
P2 & (9,0,2) & (3,0,2) & (6,0,0) \\
P3 & (2,2,2) & (2,1,1) & (0,1,1) \\
P4 & (4,3,3) & (0,0,2) & (4,3,1) \\
\bottomrule
\end{tabular}
Safe sequence: P1 $\to$ P3 $\to$ P4 $\to$ P0 $\to$ P2 (verify Need $\le$ Work each step).
\paragraph{Resource-allocation graph}
Bipartite directed graph: \textcolor{Bittersweet}{request edge} $P\to R$, \textcolor{Bittersweet}{assignment edge} $R\to P$. With single-instance resources, a \emph{cycle} $\Leftrightarrow$ deadlock; with multi-instance resources, a cycle is necessary but not sufficient.

\paragraph{POSIX semaphore reference}
\texttt{sem\_t s; sem\_init(\&s, 0, 1);} (last arg = initial count).
\texttt{sem\_wait(\&s)} (decrement; block on 0); \texttt{sem\_post(\&s)} (increment; wake one).
Named: \texttt{sem\_open("/n", O\_CREAT, 0644, 1)} — share between unrelated processes.

\paragraph{Lock-free vs wait-free}
\textcolor{Bittersweet}{Lock-free}: at least one thread always makes progress (no global stall) — usually a CAS retry loop. \textcolor{Bittersweet}{Wait-free}: \emph{every} thread completes in a bounded number of steps. Avoids deadlock and priority inversion at the cost of complex algorithms; needs strong memory barriers.
